\documentclass[nopagenr]{UWletterhead}

% Available class options (put in the square brackets above):
% * `polski` - set the language to Polish, influences the logo and prefixes
% * `english` - set the language to English, influences the logo and prefixes

% * `arimo` - set the default font to Arimo
% * `droid` - set the default font to Droid
% * `dejavu` - set the default font to DejaVu

% * `grayscale` - transform colours to grayscale, for printing

% * `autopagenr` - automatically decide whether to put numbers of pages
% * `pagenr` - always put numbers of pages
% * `nopagenr` - never put numbers of pages

% * `10pt` use 10pt font for the main body
% * `11pt` use 11pt font for the main body
% * `12pt` use 12pt font for the main body

% The defaults are: polski, droid, autopagenr, 12pt.

%%%%%%%%%%%%%%%%%%%%%%%%%%%%%%%%%%%%%%%%%%%%%%%%%%%%%%%%%%%%%%%%%%%%%%%%%%%%%%

% If you use non-standard symbols (i.e. Polish diacretic letters), you should
% define input encoding using an option like the following ones.

% Use for utf8 encoded files
%\usepackage[utf8]{inputenc}

% Use for Windows cp1250 encoded files
%\usepackage[cp1250]{inputenc}

% Use for Latin 2 encoded files
%\usepackage[latin2]{inputenc}

%%%%%%%%%%%%%%%%%%%%%%%%%%%%%%%%%%%%%%%%%%%%%%%%%%%%%%%%%%%%%%%%%%%%%%%%%%%%%%

% Control the space between the header and the body text. Identification
% Codebook does not specify that. The command is optional.

\renewcommand{\UWheadsep}{10mm} % - 10mm by default

%%%%%%%%%%%%%%%%%%%%%%%%%%%%%%%%%%%%%%%%%%%%%%%%%%%%%%%%%%%%%%%%%%%%%%%%%%%%%%

% Control the margins of the body text. Identification
% Codebook does not specify that. Top margin influences the pages that do
% not start with the UW header. The command is optional. See
% UWletterhead.cls for a discussion on reasonable margin sizes. The values
% below are the default ones.

% Modifying these values (within reasonable ranges, as defined in
% UWletterhead.cls) does not move the elements of the layout that are
% positioned accordingly to Identification Codebook.

% The current default sets the left margin in the vertical line passing
% through the middle of the logo (eagle). Another reasonable possibility
% is to set the left margin to 56mm, in that case the left margin of the
% body text is aligned to the faculty, institute, etc names.

\geometry{
	top=30mm,
	left=35mm,
	right=30mm,
	bottom=30mm
}

%%%%%%%%%%%%%%%%%%%%%%%%%%%%%%%%%%%%%%%%%%%%%%%%%%%%%%%%%%%%%%%%%%%%%%%%%%%%%%

% Control the standard parts of the text (all of these commands are optional).

% There is an automatic system that shows only lines for non-empty values.
% To make it work correctly, if we want the given field to be empty, we should
% define it as an empty string, i.e. \renewcommand{\UWfaculty}{}
% A convoluted definition like:
% \newcommand{\myEmptyString}{}
% \renewcommand{\UWfaculty}{\myEmptyString}
% will not work.
% The same applies to the values \UWaddress, \UWphone, and \UWemail.

%\renewcommand{\UWfaculty}{} % - filled to MIM automatically
\renewcommand{\UWinstitute}{Instytut Informatyki} % - empty by default
%\renewcommand{\UWfunction}{} % - empty by default
%\renewcommand{\UWperson}{} % - empty by default

%%%%%%%%%%%%%%%%%%%%%%%%%%%%%%%%%%%%%%%%%%%%%%%%%%%%%%%%%%%%%%%%%%%%%%%%%%%%%%

% Control the address line prefixes (all of these commands are optional):

%\renewcommand{\UWprefAddress}{} % - empty by default
\renewcommand{\UWprefPhone}{telefon: } % - empty by default
%\renewcommand{\UWprefEMail}{e-mail: } % - filled by default

%%%%%%%%%%%%%%%%%%%%%%%%%%%%%%%%%%%%%%%%%%%%%%%%%%%%%%%%%%%%%%%%%%%%%%%%%%%%%%

% Control the address line values (all of these commands are optional):

%\renewcommand{\UWaddress}{ul. Banacha~2, 02-097 Warszawa, Polska}
	% - filled by default
\renewcommand{\UWphone}{+48 22 55 44 460} % - empty by default
\renewcommand{\UWemail}{iinf@mimuw.edu.pl} % - empty by default

%%%%%%%%%%%%%%%%%%%%%%%%%%%%%%%%%%%%%%%%%%%%%%%%%%%%%%%%%%%%%%%%%%%%%%%%%%%%%%

% If you really need, you can overwrite the whole footer address text.
% In that case all the above values \UWaddress, \UWprefAddress, etc
% have no effect. The default \UWfootText takes care of empty values
% and lets the phone and the email be hyperlinks. Thus, in most cases
% do not use this command.

\renewcommand{\UWfootText}{
\UWprefAddress\UWaddress\\
\UWprefPhone\UWphone\\
\UWprefEMail\UWemail
}

%%%%%%%%%%%%%%%%%%%%%%%%%%%%%%%%%%%%%%%%%%%%%%%%%%%%%%%%%%%%%%%%%%%%%%%%%%%%%%

% Begin the body of the document.
\begin{document}

% Start with \UWheader to put the header at the beginning of the first page.

\UWheader

% Use \UWdate to put the date line.

\UWdate{Warszawa, 15 kwietnia 2017}

% Use \UWaddressee to put the addressee part of the text. The first
% optional argument controls the horizontal space to the left of the
% addressee paragraph.
\UWaddressee[7.5cm]{Sz. P. Anna Kowalska\\
Kierownik Katedry\\
Znaczna Persona}

% Put your body text here.
Szanowna Pani Profesor,

%%%%%%%%%%%%%%%%%%%%%%%%%%%%%%%%%%%%%%%%%%%%%%%%%%%%%%%%%%%%%%%%%%%%%%%%%%%%%%

% Identification Codebook defines that we can use the specific font (Arimo)
% in the regular or bold variant.

Nie jest dla mnie jasne co oznacza \textbf{Definicja 27} z pracy domowej.

% On the other hand, Identification Codebook does not define italic variant
% of the font. Thus, the following command, if used with arimo package option,
% leads to a warning:

%\textit{Zadanie 24} jest raczej czytelne.

%%%%%%%%%%%%%%%%%%%%%%%%%%%%%%%%%%%%%%%%%%%%%%%%%%%%%%%%%%%%%%%%%%%%%%%%%%%%%%

% It is allowed to use math mode in the documents, however the fonts used for
% math are the standard LaTeX ones (do not comply with the choice done
% in the package).

Interesuje mnie wynik obliczenia $\sin\big(\frac{\pi}{7}\big)$?

%%%%%%%%%%%%%%%%%%%%%%%%%%%%%%%%%%%%%%%%%%%%%%%%%%%%%%%%%%%%%%%%%%%%%%%%%%%%%%

% You can use the macro \UWsignature to put a standarised signature.
% The first (optional) argument of the macro controls the horizontal
% space to the left of the signature.
% The command (standard for Latex) \\[10mm] puts a vertical space
% of a given height.

\UWsignature[7.5cm]{Z wyrazami szacunku,\\[10mm]Micha{\l} Skrzypczak}

%%%%%%%%%%%%%%%%%%%%%%%%%%%%%%%%%%%%%%%%%%%%%%%%%%%%%%%%%%%%%%%%%%%%%%%%%%%%%%

% Usually, the header is only on the first page. If you want to put the header
% more times in your file, you can force a new page by \newpage
% and put there \UWheader as well.
% If you want to create multiple separate documents by one file, the class
% option `nopagenr` (used in line 1 of this file) may be helpful. It removes
% page numbers from all the pages.

\newpage
\UWheader

Szanowna Pani Magister,

%%%%%%%%%%%%%%%%%%%%%%%%%%%%%%%%%%%%%%%%%%%%%%%%%%%%%%%%%%%%%%%%%%%%%%%%%%%%%%

% If you need to use the standarised colours, you can do it. The list of
% available colours is the following (with their RGB values):
% * UWdBlue : {  0,  68, 124}
% * UWlBlue : {  0, 146, 206}
% * UWsBlue : {  0,  23,  69}
% * UWblack : {  0,    0,  0}
% * UWwhite : {255, 255, 255}
% * UWlGray : {177, 179, 182}
% * UWdGray : {128, 130, 133}

% If the option grayscale is used then the colours \UWdBlue, \UWlBlue,
% and \UWsBlue are turned manually into grayscale, as follows:
% * UWdBlue : { 58,  58,  58}
% * UWlBlue : {120, 120, 120}
% * UWsBlue : { 21,  21,  21}

% You can colour your text using \textcolor macro.

\textcolor{UWdBlue}{Nalegam na dotrzymanie terminu!}

%%%%%%%%%%%%%%%%%%%%%%%%%%%%%%%%%%%%%%%%%%%%%%%%%%%%%%%%%%%%%%%%%%%%%%%%%%%%%%

\UWsignature{Uszanowania,\\[5mm]Micha{\l} Skrzypczak}

\end{document}

